% !TEX TS-program = pdflatex
% !TEX encoding = UTF-8 Unicode

%\makeglossaries
\usepackage{comment}
\usepackage{notoccite}
\usepackage{pdfpages}
\usepackage{cite}
\usepackage{booktabs}
\usepackage{colortbl}
\usepackage{xcolor}

% Packages
%\usepackage{slashbox,multirow}
\PassOptionsToPackage{hyphens}{url}
\usepackage{multirow}
\usepackage{enumitem}
\usepackage{ragged2e}
\usepackage[utf8]{inputenc}
\usepackage[british]{babel}
\usepackage[pdftex]{hyperref}


%\usepackage[acronym,toc,style=tree,nonumberlist,automake]{glossaries}
\usepackage[nopostdot,style=super,nonumberlist,toc]{glossaries}
\usepackage{graphicx}
\usepackage{rotating}
\usepackage{tikz}
\usepackage{mathtools}
\usepackage{wrapfig}
\usepackage{lipsum}
\usepackage{numprint}
\usepackage[nottoc]{tocbibind}
\usepackage[top=3cm, bottom=4cm, inner=3cm, outer=3cm]{geometry}
\usepackage[nohyperlinks]{acronym}
\usepackage{setspace}
\usepackage{afterpage}
\usepackage{icomma}
\usepackage{calc}
\usepackage{indentfirst}
\usepackage{amsmath}
\usepackage{mathrsfs}
\usepackage{mathbbol}
\usepackage{bm}
\usepackage{pgfplots}
\pgfplotsset{compat=1.18}
\usepgfplotslibrary{groupplots}
\usetikzlibrary{datavisualization}
\usetikzlibrary{datavisualization.formats.functions}
\usetikzlibrary{arrows.meta, shapes.geometric, positioning, fit, backgrounds, calc, decorations.pathreplacing, shadows}
\usepackage{fancybox}
\usepackage{dcolumn}
\newcolumntype{d}{D{.}{.}{-1}}
\usepackage{float}
\usepackage[section]{placeins}   % auto \FloatBarrier at every \section

\usepackage{algorithm}
\usepackage{algorithmicx}
\usepackage[noend]{algpseudocode}
\usepackage{amsthm}
\theoremstyle{definition}
\newtheorem{theorem}{Theorem}[section]
\newtheorem{lemma}[theorem]{Lemma}
\newtheorem{proposition}[theorem]{Proposition}
\newtheorem{corollary}[theorem]{Corollary}
\newtheorem{defn}{Definition}[section]
\newtheorem{example}{Example}[chapter]

\usepackage{parskip}
\usepackage{fancyhdr}	
\usepackage{eso-pic}	
\usepackage{makeidx}
\usepackage[toc,page]{appendix}
\usepackage{subcaption}
\usepackage{bbold}
\usepackage{epigraph}

\setlength\epigraphwidth{.8\textwidth}
\setlength\epigraphrule{0pt}

\DeclarePairedDelimiter\abs{\lvert}{\rvert}%
\DeclarePairedDelimiter\norm{\lVert}{\rVert}%

% Options
\hypersetup{
	colorlinks,
	citecolor = black,
	filecolor = black,
	linkcolor = black,
	urlcolor = black,
	pdfstartview = FitH
}

% space between lines
% \linespread{1.6}
\onehalfspacing

% subchapter levels (chapter=0, section=1, subsection=2, subsubsection=3)
\setcounter{secnumdepth}{3}
\setcounter{tocdepth}{3}

% defining new commands [simplify work]
\newcommand{\eg}{\emph{e.g. }}
\newcommand{\etal}{\emph{et al. }}
\newcommand{\ie}{\emph{i.e. }}
\DeclareMathOperator{\e}{e}
\DeclareMathOperator{\dd}{d\!}
%\newglossary[slg]{symbols}{sym}{sbl}{List of Symbols}

% change some predefined names
\renewcommand{\thefootnote}{\arabic{footnote}}
\renewcommand{\acronymname}{List of Acronyms}
\renewcommand{\bibname}{References}

% define thousand and decimal separators, it is necessary the use of \numprint in order to work well
\npthousandsep{,}\npthousandthpartsep{}\npdecimalsign{.}

% margins
%\let\tmp\oddsidemargin
%\let\oddsidemargin\evensidemargin
%\let\evensidemargin\tmp
%\reversemarginpar

% Create glossaries
%\makeglossaries

% Create index
\makeindex

% Chapter title settings
\usepackage{titlesec}		
\titleformat{\chapter}[display]
  {\Huge\bfseries\filcenter}
  {{\fontsize{50pt}{1em}\vspace{-4.2ex}\selectfont \textnormal{\thechapter}}}{1ex}{}[]
  \titleformat{\section}
  {\normalfont\fontsize{14}{18}\bfseries}{\thesection}{1em}{}

% Header and footer settings (Select TWOSIDE or ONESIDE layout below)							
\pagestyle{fancy}  
\renewcommand{\chaptermark}[1]{\markboth{\thechapter.\space#1}{}} 

% Select one-sided (1) or two-sided (2) page numbering
\def\layout{2}	% Choose 1 for one-sided or 2 for two-sided layout
% Conditional expression based on the layout choice
\ifnum\layout=2	% Two-sided
    \fancyhf{}			 						
	\fancyhead[LE,RO]{\nouppercase{ \leftmark}}
	\setlength{\headheight}{32pt}
	\fancyfoot[LE,RO]{\thepage}
	\fancypagestyle{plain}{			% Redefine the plain page style
		\fancyhf{}
		\renewcommand{\headrulewidth}{0pt} 		
		\fancyfoot[LE,RO]{\thepage}}	
\else			% One-sided  	
  	\fancyhf{}					
	\fancyhead[C]{\nouppercase{ \leftmark}}
	\fancyfoot[C]{\thepage}
\fi

\newcommand{\backgroundpic}[3]{
	\put(#1,#2){
	\parbox[b][\paperheight]{\paperwidth}{
	\centering
	\includegraphics[width=\paperwidth,height=\paperheight,keepaspectratio]{#3}}}}

% Caption settings (aligned left with bold name)
\usepackage[labelfont=bf, textfont=normal,
			justification=justified,
			singlelinecheck=false]{caption}
% Disable hypcap globally: \captionof in minipages cannot anchor hyperref
% targets at the float top (since there is no float), which would otherwise
% produce a 'hypcap=true will be ignored' warning at every \captionof call.
\captionsetup{hypcap=false}
